%%%%%%%%%%%%%%%%%%%%%%%%%%%%%%%%%%%%%%%%%%%%%%%%%%%%%%%%%%%%%%%%%%%%%%%%%%%%%
%% Descr:       T2000b, Kapitel 1
%% Author:      Christoph Niederer, chrisi.niederer@gmail.com
%% $Id: kapitel1.tex,v 1.5 2017/10/06 14:02:51 vollmer Exp $
%%  -*- coding: utf-8 -*-
%%%%%%%%%%%%%%%%%%%%%%%%%%%%%%%%%%%%%%%%%%%%%%%%%%%%%%%%%%%%%%%%%%%%%%%%%%%%%%%

\chapter{Einleitung}

Im Rahmen der Studienarbeit wurde ein Smart Mirror entwickelt, der mithilfe eines Spionspiegels, eines Displays und eines Raspberry Pi dem Nutzer verschiedene Informationen anzeigen kann. Ein besonderes Ziel war die Integration einer Industriekamera, die neben der Pulsmessung auch die Erkennung der aktuellen Stimmung (Mood Detection) ermöglicht. Die erfassten Werte werden in Echtzeit auf dem Display sichtbar gemacht. Das Projekt vereint Hardwareaufbau, Bildverarbeitung sowie die Entwicklung einer grafischen Oberfläche.
Es sei darauf hingewiesen, dass das vorliegende Projekt noch nicht final abgeschlossen ist, jedoch die im Kurs vereinbarte Deadline für die Projektmanagement-Abgabe erreicht wurde. 
